\bookchapter{الفصل 9: العمل والتعلم عبر الهاتف المحمول}

\chapterintro{كنت أعتقد أيضاً أن العمل على الهاتف هو مجرد “حل مؤقت”، حتى سافرت مرة في مهمة عمل ونسيت حاسوبي المحمول، واضطررت إلى تعديل عرض تقديمي وتوقيع عقد باستخدام الهاتف فقط. بصراحة، عنق الزجاجة في العمل عبر الهاتف لا يكون عادةً في الهاتف نفسه، بل في عدم إلمامك بحدود قدراته.}

\booksection{أولاً: معالجة المستندات: التحرير والإبداع في راحة اليد}

تحرير المستندات على شاشة الهاتف يختلف فعلاً عن الكمبيوتر. أكبر مشكلة هي صغر الشاشة وأصابعنا السميكة وميل التنسيق إلى الاختلال. لكن هناك حيل تجعل التجربة أفضل بكثير:

\begin{itemize}
  \item \textbf{WPS Office}: يتميز بأفضل توطين محلي، ومتوافق تماماً مع تنسيقات Microsoft (docx/xlsx/pptx)، ومجموعة أدوات PDF المدمجة فيه غنية جداً. أحب بشكل خاص وظيفة "الملاحظات الصوتية" — فعند مراجعة العقد وعدم القدرة على الكتابة، يمكنك التحدث ببضع ملاحظات إلى الهاتف ليحوّلها تلقائياً إلى تعليقات نصية تُلصق في المستند. كما يحتوي WPS على كمية هائلة من القوالب المجانية، ما يوفر الكثير من الوقت عند إعداد السيرة الذاتية أو عروض التقديم.
  \item \textbf{Microsoft 365}: إذا كانت شركتك تستخدم Office أصلياً، فهذا هو الخيار الأفضل. توافقه مع تنسيقات نسخة سطح المكتب شبه مثالي، ولا تحدث مشكلة الإزاحة الطفيفة التي تظهر أحياناً في WPS عند فتح بعض التنسيقات المعقدة. وهو متكامل بعمق مع OneDrive، فالمستند الذي تعدّله على الكمبيوتر سيكون بأحدث إصدار عند فتحه على الهاتف.
  \item \textbf{Shimo Docs/Tencent Docs}: هذان يمثلان "التعاون عبر الإنترنت". مناسبان لعدة أشخاص يحررون مستنداً واحداً في الوقت نفسه، حيث تظهر تعديلات كل شخص فوراً. استخدمت Shimo Docs مع صديق لإعداد دليل سفر، وكنا نملأ المعلومات في الجدول معاً، والكفاءة أعلى بكثير من العمل الفردي.
\end{itemize}

\textbf{سير عمل عملي}: عادتي الحالية هي — الملفات المهمة على OneDrive، والتحرير اليومي بـ WPS، والتعاون الجماعي عبر Shimo Docs. منذ أن رفعت ملفاتي إلى السحابة، لم تعد تحدث لي حالة الإحراج من “نسيان محرك أقراص USB”.

\booksection{ثانياً: تطبيقات الملاحظات: من التدوين العفوي إلى نظام المعرفة}

كثير من الناس عندما يبدؤون في تدوين الملاحظات يريدون بناء نظام معقد جداً، ثم لا يصمدون أكثر من ثلاثة أيام. نصيحتي هي البدء خفيفاً ثم التثقيل — اكتسب أولاً عادة “التدوين”، ثم فكّر في مسألة “الإدارة”.

\begin{itemize}
  \item \textbf{التدوين العفوي — أدوات النظام كافية}: تطبيق الملاحظات في iOS سريع التشغيل جداً، ويدعم مسح المستندات والكتابة اليدوية والتسجيل الصوتي. وتطبيقات الملاحظات المدمجة في علامات Android التجارية المختلفة ليست أقل جودة، ويمكن تشغيلها بسرعة من شاشة القفل. هذه الأدوات مناسبة لتدوين الأفكار المؤقتة وقوائم التسوق ونقاط الاجتماعات. المهم هو سرعة التشغيل — بمجرد ورود الفكرة يمكنك تدوينها دون انتظار تحميل التطبيق.
  \item \textbf{القراءة لاحقاً — جمع المعلومات}: وظيفة "النافذة العائمة" في WeChat هي في الواقع أداة سحرية خفية. عندما تصادف مقالاً جيداً ولا يتسع الوقت لقراءته، انقر على "النافذة العائمة" في الزاوية العلوية اليمنى، فتنكمش إلى نقطة دائرية صغيرة تلتصق بحافة الشاشة، ويمكن فتحها في أي وقت لمتابعة القراءة. وهذا أسهل بكثير من إعادة التوجيه إلى "مساعد نقل الملفات". كما يمكن لأدوات القراءة لاحقاً مثل Cubox أو Pocket جمع المحتوى الجيد الذي تصادفه في أي مكان عبر منصات متعددة.
  \item \textbf{بناء النظام — Notion وFeishu}: عندما تشعر أن ملاحظاتك المتفرقة بدأت يصعب العثور عليها، حان وقت Notion. هذه الأداة مرنة جداً — يمكنك بناء "قاعدة بيانات لملاحظات القراءة" تُنظَّم تلقائياً حسب عنوان الكتاب والمؤلف والتقييم؛ أو بناء "لوحة مشاريع" لتصور تقدّم مهام العمل. "الجداول متعددة الأبعاد" في Feishu مشابهة أيضاً، وسرعة الوصول إليها محلياً أسرع. أستخدم Notion لإدارة جميع مواد الكتابة وملاحظات التعلم وتقدّم المشاريع منذ أكثر من ثلاث سنوات.
\end{itemize}

\textbf{استراتيجية ملاحظات عملية}: عادتي الشخصية هي “ثلاثية” — الأفكار المؤقتة في ملاحظات النظام (تشغيل فوري)، والمقالات للقراءة لاحقاً في النافذة العائمة لـ WeChat (في المتناول)، والمعرفة التي تحتاج تراكماً طويل الأمد في Notion (منهجية). لا تحاول من البداية حشر كل شيء في أداة واحدة، فهذا يسهّل الاستسلام.

\booksection{ثالثاً: المسح والتعرف الضوئي على الحروف (OCR): نقل العالم الورقي إلى الهاتف}

مسح المستندات بالهاتف، إذا أتقنته، يوفر الكثير من متاعب الذهاب إلى محل الطباعة. في المرة الأخيرة التي تقدمت فيها لطلب تأشيرة كنت بحاجة إلى نسخ إلكترونية من مجموعة من المستندات، وقد أنجزت ذلك كله بالهاتف.

\begin{itemize}
  \item \textbf{مستخدمو iPhone}: وظيفة المسح المدمجة في تطبيق الملاحظات على iOS جيدة جداً. افتح الملاحظات انقر على أيقونة الكاميرا اختر "مسح المستند"، ووجّه الكاميرا إلى الملف وسيتعرف تلقائياً على الحواف ويصحّح المنظور. بعد التصوير يمكنك ضبط الألوان — اختر وضع "أبيض وأسود" للحصول على أقرب نتيجة إلى الماسح الضوئي، و"تدرّج رمادي" للمواد التي تحتوي على صور.
  \item \textbf{مستخدمو Android}: العديد من العلامات التجارية توفر وضع "مسح المستندات" في الكاميرا المدمجة. إذا لم تجده، فنزّل Microsoft Lens من متجر التطبيقات (مجاني، ومتكامل مع OneDrive).
  \item \textbf{التعرف الضوئي على الحروف (OCR)}: إذا كنت تريد استخراج النص من المستندات الممسوحة وتحريره، فأنت بحاجة إلى OCR. وظيفة "استخراج النص" المدمجة في WeChat هي الأسهل — اضغط مطولاً على الصورة لاستخدامها، ودقة التعرف جيدة. للكميات الكبيرة، يمكن استخدام تطبيق Baimiao أو CamScanner. بالنسبة للتعرف على الحروف الصينية، فإن محرك OCR من Baidu يعطي نتائج أفضل؛ ووظيفة تحويل الصور إلى نص في WPS دقيقة جداً أيضاً.
  \item \textbf{نصيحة مسح صغيرة}: صوّر في مكان ذي إضاءة كافية، ويُفضّل أن يكون هناك مصباح جانبي لتجنب ظل يدك على المستند. إذا كان الملف مطوياً من زاوية، فافرده بيدك قبل التصوير. يمكن مسح المستندات متعددة الصفحات على التوالي ثم دمجها في ملف PDF واحد، بدلاً من حفظ كل صفحة على حدة.
\end{itemize}

\booksection{رابعاً: كيف تختار التخزين السحابي: لا يوجد الأفضل، بل الأنسب لك}

فكّرت طويًا في اختيار التخزين السحابي، واكتشفت في النهاية أن الأمر يعتمد على ما هو “نظامك الأساسي”.

\begin{itemize}
  \item \textbf{منظومة Apple الكاملة iCloud}: تزامن سلس يجعلك لا تشعر بوجوده. الصور التي تلتقطها بهاتفك ستكون على الكمبيوتر بعد دقائق. كما أن حماية الخصوصية من Apple هي الأكثر صرامة. لكن العيوب واضحة أيضاً: المشاركة عبر المنصات مع مستخدمي Android أو Windows مزعجة، والمساحة المجانية 5GB فقط، والتوسعة تكلف المال.
  \item \textbf{Windows كحاسوب رئيسي OneDrive}: إذا كنت تستخدم حاسوب Windows + Office 365، فإن OneDrive هو الشريك المثالي. عند حفظ مستند في Word، يُحفظ تلقائياً في السحابة. مساحة 1TB مع اشتراك Office صفقة جيدة. ووظيفة النسخ الاحتياطي التلقائي للصور في تطبيق الهاتف ممتازة أيضاً.
  \item \textbf{المشاركة محلياً Baidu Netdisk}: بصراحة، يُنتقد تقييد سرعة التنزيل في Baidu Netdisk دائماً، لكن منظومة المشاركة فيه لا تُضاهى. الكثير من المواد التعليمية والموارد المرئية على الإنترنت الصيني تُشارك عبر Baidu Netdisk. إذا كنت تنزّل ملفات يشاركها الآخرون من حين لآخر، فيمكنك الاستغناء عن العضوية المدفوعة. أما إذا كنت جامع موارد كثيف، فاشترك في العضوية، فالوقت الذي توفره يساوي أكثر من ثمنها.
  \item \textbf{Aliyun Drive/Quark Drive}: هذان من الطارئين الجدد في السنوات الأخيرة. Aliyun Drive لا يقيّد سرعة الرفع والتنزيل (حالياً)، وQuark Drive مرتبط بعمق مع متصفح Quark، فيمكنك حفظ المواد مباشرة إلى التخزين السحابي عند البحث، وهو مريح جداً.
\end{itemize}

\textbf{حلتي المختلطة}: iCloud للصور الشخصية والملفات اليومية (الأولوية للخصوصية)، وBaidu Netdisk للموارد واستقبال المشاركات (الأولوية للمنظومة)، وOneDrive لمستندات العمل (الأولوية لتكامل Office). ثلاثة خدمات تخزين سحابي، لكل منها دور، دون تداخل.

\booksection{خامساً: سيناريو عمل حقيقي عبر الهاتف المحمول}

دعني أروي تجربة حقيقية لي. سافرت مرة إلى مدينة أخرى، وطلب العميل فجأة تعديل بنود العقد. كان سير عملي كالتالي:

\begin{enumerate}
  \item تلقيت في الفندق ملاحظات التعديل من العميل عبر WeChat (رسالة نصية + صورة لملاحظات مكتوبة بخط اليد).
  \item فتحت OneDrive للعثور على العقد الأصلي، وفتحته بـ WPS للتحرير، وعدّلت البنود واحداً تلو الآخر بالرجوع إلى رسائل العميل على WeChat.
  \item بعد الانتهاء، صدّرت الملف إلى PDF باستخدام وظيفة "إخراج كـ PDF" في WPS، وأرسلته إلى مجموعة العميل على WeChat للتأكيد.
  \item طرح العميل جولة ثانية من الملاحظات، فاستخدمت أداة القلم في WPS لتدوير مواضع التعديل مباشرة على PDF، مع ملاحظات صوتية لشرح فكرة التعديل.
  \item بعد التأكد من عدم وجود أخطاء، استخدمت وظيفة "التوقيع الإلكتروني" في WPS للتوقيع عبر الإنترنت، وأرسلته إلى الطرف الآخر للتوقيع كذلك.
  \item أخيراً، حفظت نسخة من العقد الموقّع في Baidu Netdisk وiCloud، ودوّنت ملاحظة في سجل العمل على Notion.
\end{enumerate}

استغرقت العملية بأكملها من البداية إلى النهاية نحو أربعين دقيقة، دون تشغيل الكمبيوتر. بصراحة، قبل ثلاث سنوات لم أكن أصدق أن الهاتف قادر على فعل هذه الأشياء، لكنه أصبح قادراً عليها فعلاً. المهم هو الإلمام المسبق بهذه الأدوات، حتى لا تضطر إلى تنزيلها وتسجيل الدخول عند الحاجة.

\begin{tcbpractice}{جرّب الآن: ابنِ سير عمل المكتب المتنقل الخاص بك}
\begin{enumerate}
  \item اختر تطبيق المستندات الرئيسي (WPS أو Microsoft 365)، وسجّل الدخول واربطه بالتخزين السحابي.
  \item افتح الملاحظات أو تطبيق الملاحظات اللاصقة، واستخدم وظيفة المسح لتصوير مستند وتصديره إلى PDF وإرساله إلى نفسك.
  \item اختر تطبيق تخزين سحابي، واضبطه على النسخ الاحتياطي التلقائي لصور الهاتف (اختر "النسخ الاحتياطي عبر Wi-Fi فقط" لتوفير البيانات).
  \item أنشئ صفحة جديدة في Notion أو Feishu، وحاول بناء قالب بسيط لـ "سجل العمل".
\end{enumerate}

\end{tcbpractice}
